% Options for packages loaded elsewhere
\PassOptionsToPackage{unicode}{hyperref}
\PassOptionsToPackage{hyphens}{url}
\documentclass[
  brazilian,
  12pt,
]{article}
\usepackage{xcolor}
\usepackage{amsmath,amssymb}
\newcommand{\ganttbar}{\rule{0.55em}{1.15ex}}
\setcounter{secnumdepth}{-\maxdimen} % remove section numbering
\usepackage{iftex}
\ifPDFTeX
  \usepackage[T1]{fontenc}
  \usepackage[utf8]{inputenc}
  \usepackage{textcomp} % provide euro and other symbols
\else % if luatex or xetex
  \usepackage{unicode-math} % this also loads fontspec
  \defaultfontfeatures{Scale=MatchLowercase}
  \defaultfontfeatures[\rmfamily]{Ligatures=TeX,Scale=1}
\fi
\usepackage{lmodern}
\ifPDFTeX\else
  % xetex/luatex font selection
\fi
% Use upquote if available, for straight quotes in verbatim environments
\IfFileExists{upquote.sty}{\usepackage{upquote}}{}
\IfFileExists{microtype.sty}{% use microtype if available
  \usepackage[]{microtype}
  \UseMicrotypeSet[protrusion]{basicmath} % disable protrusion for tt fonts
}{}
\makeatletter
\@ifundefined{KOMAClassName}{% if non-KOMA class
  \IfFileExists{parskip.sty}{%
    \usepackage{parskip}
  }{% else
    \setlength{\parindent}{0pt}
    \setlength{\parskip}{6pt plus 2pt minus 1pt}}
}{% if KOMA class
  \KOMAoptions{parskip=half}}
\makeatother
\usepackage{longtable,booktabs,array}
\newcounter{none} % for unnumbered tables
\usepackage{calc} % for calculating minipage widths
% Correct order of tables after \paragraph or \subparagraph
\usepackage{etoolbox}
\makeatletter
\patchcmd\longtable{\par}{\if@noskipsec\mbox{}\fi\par}{}{}
\makeatother
% Allow footnotes in longtable head/foot
\IfFileExists{footnotehyper.sty}{\usepackage{footnotehyper}}{\usepackage{footnote}}
\makesavenoteenv{longtable}
\ifLuaTeX
\usepackage[bidi=basic,shorthands=off]{babel}
\else
\usepackage[bidi=default,shorthands=off]{babel}
\fi
\ifLuaTeX
  \usepackage{selnolig} % disable illegal ligatures
\fi
\setlength{\emergencystretch}{3em} % prevent overfull lines
\providecommand{\tightlist}{%
  \setlength{\itemsep}{0pt}\setlength{\parskip}{0pt}}
\usepackage{bookmark}
\IfFileExists{xurl.sty}{\usepackage{xurl}}{} % add URL line breaks if available
\urlstyle{same}
\hypersetup{
  pdflang={pt-BR},
  hidelinks,
  pdfcreator={LaTeX via pandoc}}

\author{}
\date{}

\begin{document}

INTELIGÊNCIA ARTIFICIAL MULTIAGENTE NO ENSINO SUPERIOR:UM GUIA PRÁTICO
PARA PESQUISA CIENTÍFICA ASSISTIDA E ÉTICA

\textbf{Linha de Pesquisa:} Linha 1: Inovações e Práticas em Tecnologia
Educacional

\textbf{Primeira Opção de Tema:} Tema 03: Aplicação de Inteligência
Artificial no Desenvolvimento de Ferramentas para Suporte ao Contexto
Educacional

\textbf{Segunda Opção de Tema:} Tema 04: Avaliação de Aspectos da LGPD e
Privacidade de Dados no Contexto Educacional

\textbf{Tipo de Produto Educacional:} Conteúdo educacional digital

\begin{center}\rule{0.5\linewidth}{0.5pt}\end{center}

\subsection{RESUMO}\label{resumo}

A disseminação de ferramentas de Inteligência Artificial (IA) generativa
no ensino superior introduziu riscos de alucinações, plágio
involuntário, vazamento de dados e perda de rastreabilidade de fontes.
Este anteprojeto propõe o desenvolvimento de um guia prático de uso
ético de uma plataforma de IA multiagente de código aberto, que coordena
125 agentes especializados com auditoria caixa branca e rastreabilidade
por Digital Object Identifier (DOI), como ferramenta de suporte à
pesquisa acadêmica assistida. A pesquisa adota metodologia mista em três
fases: (1) análise documental da arquitetura do ecossistema frente às
normativas de ética em IA (Lei Geral de Proteção de Dados Pessoais,
LGPD, Lei nº 13.709/2018; Resolução PRPPG/UFC nº 39/2025); (2)
desenvolvimento do guia prático validado por especialistas; (3) estudo
de caso com grupo focal de pesquisadores de pós-graduação. O produto
educacional será um manual digital interativo que sistematiza boas
práticas de uso ético de IA multiagente na pesquisa.

\textbf{Palavras-chave:} IA Multiagente. Ética na Pesquisa. Código
Aberto. LGPD. Tecnologia Educacional.

\begin{center}\rule{0.5\linewidth}{0.5pt}\end{center}

\subsection{JUSTIFICATIVA E DELIMITAÇÃO DO
PROBLEMA}\label{justificativa-e-delimitauxe7uxe3o-do-problema}

O avanço acelerado das IAs generativas transformou a pesquisa acadêmica,
mas expôs vulnerabilidades críticas: (a) modelos geram conteúdo
plausível, porém factualmente incorreto, incluindo referências
bibliográficas inexistentes; (b) o usuário não consegue auditar como o
modelo chegou a uma conclusão, violando o princípio da reprodutibilidade
científica; (c) pesquisadores inserem dados inéditos em plataformas que
os utilizam para treinamento, infringindo a LGPD (Lei nº 13.709/2018);
(d) a paráfrase automatizada pode reproduzir trechos protegidos sem
atribuição (FLORIDI, 2023; RUSSELL; NORVIG, 2022).

O Edital nº 01/2026 do PPGTE/UFC exige que todo candidato declare o uso
ou não uso de IA (Anexo IV). Essa exigência evidencia uma lacuna: faltam
diretrizes práticas e ferramentas validadas que permitam ao pesquisador
utilizar IA de forma ética, rastreável e em conformidade com a lei.

A plataforma de IA multiagente que fundamenta este anteprojeto,
disponível publicamente em repositório de código aberto, preenche essa
lacuna ao implementar: (1) auditoria de cada afirmação vinculada a DOIs
verificáveis; (2) logs imutáveis com hash SHA-256; (3) debate entre
agentes com 38 tipos de raciocínio e 10 estratégias de Teoria dos Jogos
(NASH, 1950); (4) corretor automático de vazamentos de dados.

O problema de pesquisa delimita-se à questão: como sistematizar, por
meio de um guia prático validado, o uso ético de uma plataforma de IA
multiagente como ferramenta de suporte à pesquisa científica no ensino
superior, em conformidade com a LGPD e as normativas de integridade
acadêmica da UFC? A relevância justifica-se nas dimensões acadêmica
(conhecimento original sobre IA multiagente e ética na pesquisa),
tecnológica (democratização de ferramenta gratuita e de código aberto) e
social (diretrizes replicáveis para uso responsável de IA, alinhadas aos
Objetivos de Desenvolvimento Sustentável (ODS) 4: Educação de
Qualidade).

\begin{center}\rule{0.5\linewidth}{0.5pt}\end{center}

\subsection{OBJETIVOS}\label{objetivos}

\subsubsection{Objetivo Geral}\label{objetivo-geral}

Desenvolver e validar um guia prático de uso ético de uma plataforma de
IA multiagente como ferramenta de suporte à pesquisa científica
assistida no ensino superior, em conformidade com a LGPD e as normativas
de integridade acadêmica da UFC.

\subsubsection{Objetivos Específicos}\label{objetivos-especuxedficos}

\begin{enumerate}
\def\labelenumi{\arabic{enumi}.}
\item
  Analisar a arquitetura multiagente da plataforma (125 agentes,
  pipeline de escrita científica, sistema de auditoria caixa branca)
  quanto à adequação aos princípios de ética, rastreabilidade e
  privacidade de dados no contexto educacional.
\item
  Mapear as normativas vigentes sobre uso de IA na pesquisa (Resolução
  PRPPG/UFC nº 39/2025, LGPD, Marco Civil da Internet) e identificar
  lacunas entre as exigências legais e as práticas atuais dos
  pesquisadores.
\item
  Sistematizar boas práticas em formato de guia digital, organizado nos
  módulos: (a) configuração ética do ambiente; (b) pesquisa
  bibliográfica com rastreabilidade DOI; (c) redação assistida com
  auditoria de citações; (d) proteção de dados sensíveis conforme a
  LGPD.
\item
  Validar o guia por meio de estudo de caso com grupo focal de 8 a 12
  pesquisadores de pós-graduação, avaliando usabilidade, efetividade na
  prevenção de más condutas e percepção de utilidade.
\item
  Avaliar os aspectos de privacidade e proteção de dados do ecossistema,
  verificando conformidade com a LGPD e propondo recomendações de
  aprimoramento (Tema 04).
\end{enumerate}

\begin{center}\rule{0.5\linewidth}{0.5pt}\end{center}

\subsection{FUNDAMENTAÇÃO TEÓRICA}\label{fundamentauxe7uxe3o-teuxf3rica}

A fundamentação estrutura-se em três eixos interdisciplinares:

\subsubsection{Eixo 1: IA e Sistemas
Multiagentes}\label{eixo-1-ia-e-sistemas-multiagentes}

Russell e Norvig (2022) definem agente inteligente como entidade que
percebe o ambiente por sensores e age por atuadores. Sistemas
multiagentes (SMA) estendem esse conceito: múltiplos agentes autônomos
cooperam e competem para resolver problemas complexos. Diferentemente
dos Large Language Models (LLMs) monolíticos, que processam
\emph{prompts} em única inferência estatística, SMAs distribuem o
raciocínio entre agentes especializados, cada um com conhecimento de
domínio e protocolos de validação específicos. Wooldridge (2009) aponta
três vantagens dos SMAs: robustez (falha de um agente não compromete o
sistema), escalabilidade e explicabilidade (rastreabilidade de cada
decisão a um agente específico); propriedades que os tornam ideais para
aplicações educacionais que exigem auditabilidade.

A plataforma que fundamenta este anteprojeto concretiza esses princípios
em um pipeline de escrita científica no qual 49 agentes colaboram em 8
estágios para produzir artigos. Seu diferencial é o fórum de debate
entre agentes, que utiliza 38 tipos de raciocínio e 10 estratégias de
Teoria dos Jogos para confrontar hipóteses e validar conclusões antes de
apresentá-las ao usuário (NASH, 1950).

\subsubsection{Eixo 2: Tecnologia Educacional e Letramento
Digital}\label{eixo-2-tecnologia-educacional-e-letramento-digital}

Valente (2014) argumenta que a integração de tecnologias digitais no
ensino superior requer \emph{transposição didática} que transforme
práticas pedagógicas, não apenas substitua ferramentas. Moran (2018)
complementa que metodologias ativas apoiadas por tecnologia deslocam o
professor para mediador e o aluno para protagonista. Siemens (2004), com
o Conectivismo, concebe a aprendizagem como construção de redes entre
nodos de informação: artigos, bases de dados, pesquisadores e agentes de
IA, valorizando a capacidade de navegar, filtrar e sintetizar
informações. O letramento digital crítico (FREIRE, 1996) oferece a lente
pedagógica para distinguir o uso instrumental do uso crítico-reflexivo
das ferramentas de IA.

\subsubsection{Eixo 3: Ética da IA, Privacidade e Conformidade
Legal}\label{eixo-3-uxe9tica-da-ia-privacidade-e-conformidade-legal}

Floridi (2023) estabelece cinco princípios para IA ética: beneficência,
não maleficência, autonomia, justiça e explicabilidade. A plataforma
alinha-se a todos: beneficência na aceleração da produção científica com
qualidade; não maleficência pelos filtros de plágio e corretor
automático; autonomia porque o sistema é assistente (não substituto) do
pesquisador; justiça pelo código aberto e gratuito; explicabilidade pela
arquitetura de auditoria caixa branca. A LGPD (BRASIL, 2018) e a
Resolução PRPPG/UFC nº 39/2025 operacionalizam esses princípios na
pesquisa acadêmica. O orientador dos Temas 03 e 04 do PPGTE atua
precisamente na interseção entre desenvolvimento de ferramentas de IA
para suporte educacional e avaliação de privacidade de dados, conferindo
aderência temática ao projeto.

\begin{center}\rule{0.5\linewidth}{0.5pt}\end{center}

\subsection{METODOLOGIA}\label{metodologia}

Pesquisa de abordagem mista (qualitativa-quantitativa), natureza
aplicada, objetivo exploratório-descritivo, estruturada em três fases:

\textbf{Fase 1: Análise Documental e Arquitetural (Meses 1-4):}
mapeamento da arquitetura da plataforma (125 agentes, pipeline de
escrita científica, sistema de auditoria); revisão documental da LGPD e
Resolução PRPPG/UFC nº 39/2025; análise de conformidade entre recursos
de auditoria e exigências legais. Produto: Relatório técnico de
conformidade LGPD.

\textbf{Fase 2: Desenvolvimento do Guia Prático (Meses 5-12):}
estruturação em 4 módulos: (A) configuração ética do ambiente; (B)
pesquisa bibliográfica com rastreabilidade DOI; (C) redação acadêmica
com auditoria integrada; (D) proteção de dados sensíveis e protocolos de
anonimização. Validação por painel de 3 especialistas (IA, Direito
Digital/LGPD, Educação). Produto: manual digital interativo (web
responsivo, com exemplos e checklists).

\textbf{Fase 3: Estudo de Caso com Grupo Focal (Meses 13-20):}
participantes: 8 a 12 pesquisadores de pós-graduação da UFC. Quatro
encontros quinzenais de 2h utilizando a plataforma com o guia prático.
Coleta: gravações (com consentimento), questionários Likert
pré/pós-intervenção e análise de logs. Análise: estatística descritiva e
análise de conteúdo temática (BARDIN, 2016).

\textbf{Fase 4: Sistematização e Defesa (Meses 21-24):} consolidação dos
resultados, redação da dissertação, publicação do produto educacional e
defesa pública.

\textbf{Aspectos Éticos:} projeto submetido ao Comitê de Ética em
Pesquisa (CEP) da UFC; todos os participantes assinarão Termo de
Consentimento Livre e Esclarecido (TCLE); nenhum dado sensível será
inserido em plataformas externas; processamento local na plataforma
multiagente.

\begin{center}\rule{0.5\linewidth}{0.5pt}\end{center}

\subsection{CRONOGRAMA}\label{cronograma}

{\def\LTcaptype{none} % do not increment counter
\begin{longtable}[]{@{}lcccc@{}}
\toprule\noalign{}
Atividade & Sem. 1 & Sem. 2 & Sem. 3 & Sem. 4 \\
\midrule\noalign{}
\endhead
\bottomrule\noalign{}
\endlastfoot
Revisão de literatura & \ganttbar{}\ganttbar{}\ganttbar{}\ganttbar{}\ganttbar{}\ganttbar{}\ganttbar{}\ganttbar{} & \ganttbar{}\ganttbar{}\ganttbar{}\ganttbar{} & & \\
Fase 1: Análise documental e arquitetural & \ganttbar{}\ganttbar{}\ganttbar{}\ganttbar{}\ganttbar{}\ganttbar{}\ganttbar{}\ganttbar{} & & & \\
Fase 2: Desenvolvimento do guia prático & & \ganttbar{}\ganttbar{}\ganttbar{}\ganttbar{}\ganttbar{}\ganttbar{}\ganttbar{}\ganttbar{} & \ganttbar{}\ganttbar{}\ganttbar{}\ganttbar{}\ganttbar{}\ganttbar{} & \\
Validação por especialistas & & & \ganttbar{}\ganttbar{}\ganttbar{}\ganttbar{} & \\
Fase 3: Estudo de caso (grupo focal) & & & \ganttbar{}\ganttbar{}\ganttbar{}\ganttbar{}\ganttbar{}\ganttbar{}\ganttbar{}\ganttbar{} & \ganttbar{}\ganttbar{}\ganttbar{}\ganttbar{}\ganttbar{}\ganttbar{} \\
Análise dos dados & & & & \ganttbar{}\ganttbar{}\ganttbar{}\ganttbar{}\ganttbar{}\ganttbar{} \\
Redação da dissertação & & \ganttbar{}\ganttbar{} & \ganttbar{}\ganttbar{}\ganttbar{}\ganttbar{} & \ganttbar{}\ganttbar{}\ganttbar{}\ganttbar{}\ganttbar{}\ganttbar{}\ganttbar{}\ganttbar{} \\
Publicação do produto educacional & & & & \ganttbar{}\ganttbar{}\ganttbar{}\ganttbar{}\ganttbar{}\ganttbar{} \\
Defesa pública & & & & \ganttbar{}\ganttbar{}\ganttbar{}\ganttbar{} \\
\end{longtable}
}

\begin{center}\rule{0.5\linewidth}{0.5pt}\end{center}

\subsection{REFERÊNCIAS}\label{referuxeancias}

BARDIN, L. \textbf{Análise de conteúdo}. São Paulo: Edições 70, 2016.
ISBN 978-85-62938-04-7.

BRASIL. \textbf{Lei nº 13.709, de 14 de agosto de 2018}. Lei Geral de
Proteção de Dados Pessoais (LGPD). Diário Oficial da União, Brasília,
DF, 15 ago. 2018. Disponível em:
https://www.planalto.gov.br/ccivil\_03/\_ato2015-2018/2018/lei/l13709.htm.
Acesso em: 24 maio 2026.

FLORIDI, L. \textbf{The Ethics of Artificial Intelligence}: principles,
challenges, and opportunities. Oxford: Oxford University Press, 2023.
DOI: https://doi.org/10.1093/oso/9780198883098.001.0001.

FREIRE, P. \textbf{Pedagogia da autonomia}: saberes necessários à
prática educativa. São Paulo: Paz e Terra, 1996. ISBN 978-85-7753-015-1.

MORAN, J. M. \textbf{Metodologias ativas para uma educação inovadora}:
uma abordagem teórico-prática. Porto Alegre: Penso, 2018. ISBN
978-85-8429-116-8.

NASH, J. F. Equilibrium points in n-person games. \textbf{Proceedings of
the National Academy of Sciences}, v. 36, n.~1, p.~48-49, 1950. DOI:
https://doi.org/10.1073/pnas.36.1.48.

RUSSELL, S.; NORVIG, P. \textbf{Inteligência artificial}: uma abordagem
moderna. 4. ed.~Rio de Janeiro: GEN LTC, 2022. ISBN 978-85-216-3749-3.

SIEMENS, G. Conectivismo: uma teoria de aprendizagem para a era digital.
\textbf{International Journal of Instructional Technology and Distance
Learning}, v. 2, n.~1, p.~3-10, 2004. Disponível em:
http://www.itdl.org/Journal/Jan\_05/article01.htm. Acesso em: 24 maio
2026.

UNIVERSIDADE FEDERAL DO CEARÁ. Pró-Reitoria de Pesquisa e Pós-Graduação.
\textbf{Resolução PRPPG/UFC nº 39, de 1º de outubro de 2025}. Dispõe
sobre o uso de Inteligência Artificial na pesquisa acadêmica. Fortaleza:
UFC, 2025. Disponível em: https://ppgte.ufc.br/pt/editais/. Acesso em:
24 maio 2026.

UNIVERSIDADE FEDERAL DO CEARÁ. Programa de Pós-Graduação em Tecnologia
Educacional. \textbf{Edital nº 01/2026}: processo seletivo ao Mestrado
Profissional em Tecnologia Educacional. Fortaleza: PPGTE/UFC, 2026.
Disponível em: https://ppgte.ufc.br/pt/editais/. Acesso em: 24 maio
2026.

VALENTE, J. A. A comunicação e a educação baseada no uso das tecnologias
digitais de informação e comunicação. \textbf{UNOPAR Científica Ciências
Humanas e Educação}, Londrina, v. 1, n.~1, p.~141-166, 2014.

WOOLDRIDGE, M. \textbf{An Introduction to MultiAgent Systems}. 2.
ed.~Chichester: John Wiley \& Sons, 2009. ISBN 978-0-470-51946-2.



% ======================================================================
% APÊNDICE: Validação TDD das SPECs (Meta-Avaliação DCA)
% ======================================================================
\newpage
\section*{Apêndice: Validação TDD das SPECs da Meta-Avaliação DCA}
\addcontentsline{toc}{section}{Apêndice: Validação TDD das SPECs}

O sistema \textbf{TDDAcademic v2.0} implementa validação automatizada das 5 Especificações Técnicas extraídas da Meta-Avaliação DCA, utilizando desenvolvimento orientado a testes (TDD) com arquitetura orientada a objetos.

\subsection*{Arquitetura do Sistema}
\addcontentsline{toc}{subsection}{Arquitetura do Sistema}

\begin{itemize}
    \item \textbf{SpecValidator} (ABC): Classe abstrata com Template Method \texttt{validar()}
    \item \textbf{PciValidator}: Calibração do Process Confidence Index (SPEC-PCI-001)
    \item \textbf{CodeValidator}: Verificador de obrigatoriedade de código (SPEC-CODE-001)
    \item \textbf{AntisymValidator}: Verificador de antisimetria do produto exterior (SPEC-ANTISYM-001)
    \item \textbf{NarrValidator}: Conversor de narração para demonstração (SPEC-NARR-001)
    \item \textbf{CoraValidator}: Expansão do escopo do Cora-Debate (SPEC-CORA-001)
    \item \textbf{TestRunner}: Orquestrador com Factory Method
    \item \textbf{TestReport}: Relatório com saída JSON e Markdown
\end{itemize}

Design Patterns: Template Method, Factory Method, Strategy, Composite, Value Object.

\subsection*{Resultados da Execução}
\addcontentsline{toc}{subsection}{Resultados da Execução}

\begin{longtable}[]{@{}lllc@{}}
\toprule
SPEC & Descrição & Testes & Taxa \\
\midrule
\endhead
PCI-001 & Calibração do PCI & 5/5 & 100\% \\
CODE-001 & Obrigatoriedade de Código & 5/5 & 100\% \\
ANTISYM-001 & Antisimetria do Produto Exterior & 5/5 & 100\% \\
NARR-001 & Narração $\rightarrow$ Demonstração & 5/5 & 100\% \\
CORA-001 & Expansão do Cora-Debate & 5/5 & 100\% \\
\midrule
\textbf{Total} & & \textbf{25/25} & \textbf{100\%} \\
\bottomrule
\end{longtable}

\subsection*{Lições Aprendidas}
\addcontentsline{toc}{subsection}{Lições Aprendidas}

\begin{enumerate}
    \item \textbf{Regex flexível:} Comutadores LaTeX aceitam \texttt{{[}{]}}, \texttt{\{\}}, \texttt{\textbackslash{[}\textbackslash{]}} com/sem escaping
    \item \textbf{Normalização:} Subscritos \texttt{J\_1} e \texttt{J1} são normalizados antes da comparação
    \item \textbf{Pullback Hénon:} Erro clássico (+b em vez de -b) detectado por contexto $F^{*}$
    \item \textbf{Falsos positivos:} Texto matemático puro (Stokes, Laplace) não dispara CODE-001
    \item \textbf{Sinal vs. magnitude:} Violação de sinal é mais comum que violação de magnitude em álgebras de Lie
\end{enumerate}

\subsection*{Como Executar}
\addcontentsline{toc}{subsection}{Como Executar}

\begin{verbatim}
python tdd_academic_validator.py
\end{verbatim}

Arquivos gerados: \texttt{relatorio\_tdd\_specs.json}, \texttt{relatorio\_tdd\_specs.md}.


% ======================================================================
% APÊNDICE B: Análise do Mapa Taxonômico (150 Questões)
% ======================================================================
\newpage
\section*{Apêndice B: Análise do Mapa Taxonômico --- 150 Questões}
\addcontentsline{toc}{section}{Apêndice B: Análise do Mapa Taxonômico}

Este apêndice apresenta o mapeamento taxonômico completo de 150 questões
(60 EUF + 90 ENA) segundo a Taxonomia 350 de Raciocínios Científicos e a
Taxonomia de Bloom. A análise inclui:
\begin{itemize}
  \item Classificação de 32 tipos de raciocínio (R-codes) em 5 categorias
    taxonômicas (Lógico-Dedutivos, Analítico-Estruturais,
    Econômico-Decisórios, Físico-Matemáticos e Lógica Matemática
    Avançada);
  \item Distribuição de Bloom (Aplicar/Analisar) por exame, por área
    e por R-code;
  \item Correlações estatísticas ($\chi^2$ e V de Cramer) entre área,
    R-code e nível de Bloom;
  \item Comparação entre os perfis taxonômicos da EUF (Física) e do
    ENA (Matemática).
\end{itemize}

O relatório completo em PDF encontra-se no arquivo:
\texttt{dados\_entrada/mestrado/analise\_mapa\_taxonomico.pdf}

\vspace{0.3cm}
\noindent\textbf{Resultados Principais:}
\begin{itemize}
  \item V de Cramer Área $\times$ R-code = 0,770 (associação muito forte);
  \item V de Cramer Bloom $\times$ R-code = 0,602 (associação forte);
  \item V de Cramer Área $\times$ Bloom = 0,232 (associação fraca);
  \item EUF: domínio da Cat.~XIII (Físico-Matemáticos, 86,7\%);
  \item ENA: domínio da Cat.~I (Lógico-Dedutivos, 67,8\%);
  \item Bloom total (150 questões): Aplicar = 87 (58\%), Analisar = 63 (42\%).
\end{itemize}

\vspace{0.3cm}
\noindent\textbf{Arquivos:}
\begin{itemize}
  \item \texttt{analise\_mapa\_taxonomico.pdf}: relatório completo (15 páginas)
  \item \texttt{analise\_mapa\_taxonomico.tex}: código-fonte LaTeX
  \item \texttt{dataset\_completo.json}: dados de 150 questões
  \item \texttt{comparativo\_euf.html}: dashboard interativo (60 EUF)
\end{itemize}

\end{document}

